\chapter{Testing, Deployment, and Evaluation}
\label{ch:testing}

This chapter describes how \voname{} is tested and run, and then evaluates the
system against the requirements of Chapter~\ref{ch:requirements}. In keeping with
the rest of the book, the evaluation is honest: it records both what works
end-to-end and what remains partial or unimplemented, so the account is a faithful
picture of a substantial but still-evolving system.

\section{Testing Strategy}

\voname{} is tested at several levels, though not uniformly across services.

\begin{description}
  \item[Unit tests] exercise individual components in isolation --- service
  methods, mappers, and validation logic. The workspace service is the most
  rigorously tested: its build enforces a 90\% line-coverage threshold through
  JaCoCo, so its business logic (role checks, invitation lifecycle, optimistic
  layout updates) is well covered.
  \item[Integration tests] verify a service together with its real
  infrastructure. The Mongo/Redis-backed services (chat, room, notifications)
  use \textbf{Testcontainers} to spin up genuine MongoDB, Redis, and RabbitMQ
  instances inside the test run, so that database queries, index constraints,
  Redis scripts, and message-broker round-trips are exercised against the real
  technologies rather than mocks. The workspace service uses Failsafe-driven
  integration tests against PostgreSQL with Flyway-migrated schemas.
  \item[Contract documentation] takes the place of some cross-service testing:
  each implemented service publishes an OpenAPI specification (rendered through a
  Swagger UI in the workspace service), which fixes the request/response contract
  that clients and other services code against.
  \item[Manual and exploratory testing] covers the real-time and multi-client
  behaviours that are hardest to automate --- proximity calls, presence updates,
  notification delivery across app states --- and is supported by the
  demonstration stack described next.
\end{description}

\section{Deployment and the Demonstration Stack}

Every service ships with the container definitions for its own dependencies, and
Spring Boot's Docker Compose integration allows a service to bring those
dependencies up automatically on start. Service locations are externalised as
environment variables, so the same artefacts run unchanged from a developer laptop
to a container cluster; in a clustered deployment the gateway is the only publicly
exposed process, with everything else on an internal network.

The clearest evidence that the platform runs as an integrated whole is the
\textbf{one-command demonstration stack}. A single script:

\begin{enumerate}
  \item starts the shared infrastructure (PostgreSQL, MongoDB, Redis, RabbitMQ)
  in Docker;
  \item launches the workspace, chat, room, and gateway services, together with
  the Colyseus fork and the SkyOffice 2D client;
  \item seeds \emph{two workspaces of five members each}, consistently across
  every service --- the workspace service seeds (via a Flyway migration) a tile
  map with a meeting-room zone, five desks near the spawn point, and computers and
  whiteboards; the chat service seeds the canonical channel with those members and
  some messages; the room service seeds a shared voice room;
  \item mints a demonstration JWT and prints a ready-to-open URL for each of two
  users.
\end{enumerate}

Opening the two URLs in separate browser sessions places two avatars in the same
seeded office, close enough to fall into one proximity voice group, sharing a
channel and a voice room --- so a reviewer can, in a single command, walk two
avatars together to trigger proximity voice and exchange chat, all backed by the
real microservices. All seeders are idempotent, so the stack can be restarted
freely. This stack is both a testing aid and the canonical way to demonstrate the
system.

\section{Observability: Metrics and Dashboards}
\label{sec:observability}

Running a distributed system is only manageable if its runtime behaviour is
visible. Beyond the per-service health endpoints and application logs, \voname{}
exposes a \textbf{metrics} pipeline built on the standard Spring Boot stack and
visualised in Grafana. It covers the five implemented services --- \svc{user}
(\code{:8091}), \svc{workspace} (\code{:8087}), \svc{chat} (\code{:8084}),
\svc{room} (\code{:8086}), and \svc{notifications} (\code{:8082}).

\subsection{The Metrics Pipeline}

Every service already depends on Spring Boot Actuator; adding a Micrometer
Prometheus registry turns each one into a metrics source. Micrometer instruments
the application automatically --- HTTP server timings, JVM and process internals,
and the datastore and connection-pool clients --- and Actuator publishes the
result as a Prometheus-format text endpoint at \route{/actuator/prometheus}. A
single \textbf{Prometheus} server periodically \emph{scrapes} that endpoint on
each service and stores the samples as time series; \textbf{Grafana} then queries
Prometheus with PromQL and renders the dashboards. The flow is pull-based and
one-directional, as shown in Figure~\ref{fig:observability}.

\begin{figure}[htbp]
\centering
\begin{tikzpicture}[font=\footnotesize,
  svc/.style={draw=voPrimary,fill=voPrimary!8,rounded corners=3pt,align=center,
              minimum height=0.68cm,text width=3.0cm,inner sep=3pt},
  prom/.style={draw=voAccent!70!black,fill=voAccent!12,rounded corners=3pt,align=center,
               minimum height=1.4cm,text width=2.3cm},
  graf/.style={draw=voIndigo,fill=voIndigo!10,rounded corners=3pt,align=center,
               minimum height=1.4cm,text width=2.7cm},
  scr/.style={-{Stealth[]},voGray,thick},
]
\node[svc] (s1) at (0,3.6) {\svc{user}\, \code{:8091}};
\node[svc] (s2) at (0,2.7) {\svc{workspace}\, \code{:8087}};
\node[svc] (s3) at (0,1.8) {\svc{chat}\, \code{:8084}};
\node[svc] (s4) at (0,0.9) {\svc{room}\, \code{:8086}};
\node[svc] (s5) at (0,0.0) {\svc{notifications}\, \code{:8082}};
\node[prom] (prom) at (5.6,1.8) {\textbf{Prometheus}\\[2pt]\scriptsize scrape \&\\\scriptsize store series};
\node[graf] (graf) at (9.4,1.8) {\textbf{Grafana}\\[2pt]\scriptsize PromQL queries\\\scriptsize HTTP \& datastore\\\scriptsize dashboards};
\foreach \s in {s1,s2,s3,s4,s5}{ \draw[scr] (\s.east) to[out=0,in=180] (prom.west); }
\node[voGray,font=\scriptsize,align=center] at (3.35,3.95) {scrape (pull)\\\route{/actuator/prometheus}};
\draw[scr] (prom) -- node[above,font=\scriptsize,voGray]{query} (graf);
\end{tikzpicture}
\caption{The metrics pipeline: each implemented service exposes a Prometheus
endpoint; one Prometheus server scrapes all five and stores the series, and
Grafana queries it to render the HTTP and datastore dashboards.}
\label{fig:observability}
\end{figure}

Exposing the endpoint is a small, uniform change per service: the
\code{micrometer-registry-prometheus} dependency on the classpath, and the
endpoint added to the Actuator exposure list.

\begin{lstlisting}[caption={Enabling the Prometheus endpoint in each service (\code{application.yml}).}]
management:
  endpoints:
    web:
      exposure:
        include: health, info, prometheus
  metrics:
    tags:
      application: ${spring.application.name}
\end{lstlisting}

Prometheus is then pointed at the five services with a single scrape job.

\begin{lstlisting}[caption={The Prometheus scrape job targeting the five services (\code{prometheus.yml}).}]
scrape_configs:
  - job_name: virtual-office
    metrics_path: /actuator/prometheus
    static_configs:
      - targets: ["user:8091", "workspace:8087", "chat:8084",
                  "room:8086", "notifications:8082"]
\end{lstlisting}

Because scraping is pull-based, the services need know nothing about Prometheus
--- they simply expose the endpoint, so adding or removing a target is a change to
Prometheus's configuration alone.

\subsection{What the Dashboards Show}

The Grafana dashboards concentrate on the two signals that matter most for this
system's health: request behaviour at the edge of each service, and the state of
the datastores behind them.

\begin{description}
  \item[HTTP request rate, latency, and errors] are derived from Micrometer's
  \code{http.server.requests} timer, which records every handled request tagged
  by \code{method}, \code{uri}, and response \code{status}. From it the dashboards
  plot per-service throughput, latency percentiles (the timer captures a
  distribution, so p95/p99 are available), and the error rate as the share of
  \code{4xx}/\code{5xx} responses --- the quickest way to see a service slowing
  down or failing a dependency.
  \item[Datastore health] is read from the persistence clients each service
  already uses. The two SQL services surface \textbf{HikariCP} connection-pool
  gauges --- active, idle, and pending connections and acquisition time --- for
  \svc{user}'s MySQL and \svc{workspace}'s PostgreSQL, making pool exhaustion or
  slow queries visible before they surface as request errors; the Mongo/Redis
  services (\svc{chat}, \svc{room}, \svc{notifications}) surface their MongoDB and
  Redis client metrics on the equivalent panels.
\end{description}

Because every service is scraped uniformly, a dashboard can be read per-service or
aggregated across the fleet, so a reviewer can watch the whole backend from one
screen while exercising the demonstration stack of the previous section.

\begin{notebox}
\textbf{Scope.} Observability here is metrics-and-dashboards for the five
implemented services; it deliberately does not yet include the API gateway or the
unimplemented \svc{calendar}/\svc{tasks} scaffolds, and it stops short of
distributed tracing and alerting. The Prometheus endpoint also carries the
standard JVM and process metrics even though the current dashboards focus on HTTP
and datastore health, so those panels can be added without any change to the
services.
\end{notebox}

\section{Evaluation against Requirements}

Table~\ref{tab:eval} assesses the functional requirement groups of
Section~\ref{sec:fr}. The status column distinguishes fully working end-to-end
capabilities from those that are partial (implemented on some clients or with
known gaps) or unimplemented.

\begin{table}[htbp]
\centering
\caption{Evaluation of functional requirement groups.}
\label{tab:eval}
\small
\begin{tabularx}{\textwidth}{L{4.4cm} L{2.4cm} X}
\toprule
\textbf{Requirement group} & \textbf{Status} & \textbf{Notes} \\
\midrule
Identity \& accounts (FR-01--07) & Complete & Registration, login, JWT, OTP, profile, photo; events emitted. \\
Workspaces, desks, teams (FR-08--15) & Complete & Full PostgreSQL/Flyway service with roles, invitations, map layout, and the internal session API. \\
Messaging (FR-16--24) & Complete & Channels, DMs, threads, edits, mentions, read receipts, idempotency, event-driven channel provisioning; on both web and mobile. \\
Meeting rooms \& proximity voice (FR-25--28) & Complete (backend) & Rooms, presence, Agora tokens, and the deterministic proximity-voice engine all implemented; consumed by the web client. \\
2D office (FR-29--33) & Complete (client) & Movement, avatars, interactive objects, proximity media, in-world chat; integrated in the demo stack. \\
Notifications (FR-34--37) & Complete & E-mail and in-app paths, real-time push, idempotency. The \code{TASK\_ASSIGNED} consumer awaits its producer. \\
Web client (FR-38) & Complete & Full collaboration surface; some prototype-grade UX. \\
Mobile client (FR-39) & Partial & Chat domain implemented to a high standard; no office/workspace admin/calling. \\
Real-time auth (FR-40) & Complete & Ticket handshake and auto-reconnect on all clients. \\
Tasks \& calendar & Not implemented & Scaffolds only; the notification consumer for tasks exists. \\
\bottomrule
\end{tabularx}
\end{table}

\section{Non-functional Evaluation}

Against the non-functional targets of Section~\ref{sec:nfr}, the design choices
documented throughout the book meet the important ones:

\begin{itemize}
  \item \textbf{Performance (NFR-01)} is addressed structurally: avatar movement
  uses client-side prediction with remote interpolation, and Colyseus transmits
  only changed fields; chat delivery is push-based over persistent sockets.
  \item \textbf{Scalability (NFR-02)} is met by keeping media off the application
  servers entirely --- peer-to-peer in the office and an external SFU for rooms ---
  so server load does not grow with audio/video volume.
  \item \textbf{Security (NFR-03)} is enforced at the edge: the gateway validates
  the JWT once, strips client-supplied identity headers, and blocks internal
  endpoints, while internal calls carry a separate shared token.
  \item \textbf{Reliability and consistency (NFR-04, 05)} come from durable,
  idempotent, retried events and the fail-closed internal calls, giving eventual
  consistency that clients tolerate.
  \item \textbf{Maintainability and portability (NFR-06, 07)} follow from
  per-service data ownership and the container definitions and demo stack.
\end{itemize}

The principal non-functional gaps are operational hardening rather than design
flaws: secrets (the JWT signing key, mail credentials, the Agora certificate, and
the internal token) are currently committed as configuration defaults and should
move to a secret manager; platform-level role-based access control is not yet
modelled (the gateway assigns every user the same platform role, with real
authority resolved per workspace); and observability, while it now includes
Prometheus metrics and Grafana dashboards for the five implemented services
(Section~\ref{sec:observability}), does not yet extend to the gateway,
distributed tracing, or alerting.

\section{Limitations}

Consolidating the honest notes made throughout the book, the current limitations
are:

\begin{enumerate}
  \item The \textbf{tasks} and \textbf{calendar} services are unimplemented
  scaffolds, and the \textbf{shared library} is an empty placeholder.
  \item The \textbf{mobile} client covers only chat.
  \item The \textbf{2D office} runs as an integrated demo but is not yet driven
  fully live from the workspace service in steady-state operation.
  \item Several \textbf{secrets} are configuration defaults rather than injected
  at deploy time.
  \item Some \textbf{web UX} is prototype-grade, and a few features (presence
  roster, typing display, attachment upload) are partially wired.
\end{enumerate}

None of these undermines the core result: an integrated, multi-client, real-time
collaboration platform whose central capabilities --- identity, workspaces,
messaging, rooms, proximity communication, and notifications --- work together
end-to-end. The next chapter reflects on that result and lays out the roadmap
these limitations imply.
